% !TEX root = main.tex
\section*{3. Introduction의 순서를 어떻게 바꾸면 좋은가?}
아래는 수정한 서론의 완성 문장이 아니라, 각 section에 어떤 설명을 넣을지 정리한 순서입니다.

\subsection*{1.1 DPP에 필요한 정보와 공개하기 어려운 정보}
먼저 DPP의 환경정보를 만들려면 앞선 공급자가 제공한 자료와 생산 과정의 계산이 필요하다고 설명합니다.
자료가 올바르게 발급됐는지와 그 자료를 제품별로 나누어 계산한 결과가 맞는지는 별도로 확인해야 합니다.

이어서 기업이 상세 수치와 원료·제품 사이의 연결을 모두 공개하기 어려운 이유를 설명합니다.
이 section은 \textbf{정보를 직접 공개하지 않아도 거래가 규칙에 맞는지는 검증할 수 있지만, 그 결과만으로 실제 생산 이력을 알 수는 없다}는 지점에서 끝냅니다.

\subsection*{1.2 비공개 공급망에서도 필요한 양방향 추적}
완제품에서 출발해 원료와 공정을 조사하는 상황을 먼저 제시합니다.
그다음 문제가 있는 원료에서 출발해 이후 어떤 기록으로 이어졌는지 찾는 상황을 설명합니다.

Forward Tracing에서는 소유권이 넘어간 뒤의 소비를 연결해야 한다는 어려움을 짚습니다.
Sender에게 Receiver의 다음 소비를 알려주지 않으면서 감사자의 추적은 유지해야 한다는 문제로 연결합니다.

기존의 사후 오류 사례도 이 section에 남깁니다.
다만 Freeze부터 설명하기보다, \textbf{추적으로 대상을 찾고 사용 중단이 승인됐을 때 다음 소비를 막는다}는 순서로 설명합니다.

\subsection*{1.3 기존 연구와 이번 논문이 다룰 과제}
기존 연구가 이미 제공하는 비공개 검증, 공급망 이력, 감사와 동결 기능을 간략히 소개합니다.
그 기능들이 없다고 주장하기보다, 이번 논문에서 함께 다룰 조건을 밝힙니다.

최종 Claim에 사용한 기록, Auditor가 추적하는 기록, 원장이 소비를 허용하거나 거부하는 기록이 같아야 합니다.
여기에 Sender와 Receiver의 비밀을 구분하는 소유권 이전 방식을 함께 설명할 과제로 둡니다.

이 section의 끝에서는 \textbf{기존 기술의 이름이 아니라, 그 기술들을 사용할 때 어떤 정보와 권한을 일치시켜야 하는지}가 남아야 합니다.

\subsection*{1.4 zkDPP가 선택한 해결 방법}
Note가 공급망 Event를 거쳐 Claim으로 이어지는 과정을 먼저 보여줍니다.
그다음 Voucher가 Recall과 Receiver Note의 후속 소비 비공개성을 어떻게 지원하는지 설명합니다.

각 거래에서 남기는 감사 정보가 실제 입력·출력과 일치하도록 증명한다는 점도 연결합니다.
승인된 Auditor가 이를 복호화해 양방향으로 추적하고, 필요한 경우 미소비 대상을 동결한다고 설명합니다.

여기서는 전체 동작을 이해하는 데 필요한 내용만 남깁니다.
키 복구의 세부 단계, 소비와 동결이 경합할 때의 처리, 전체 완료 판정의 자세한 조건은 본문에서 설명합니다.

\subsection*{1.5 제안한 내용과 확인한 결과}
공급망 동작과 Claim을 연결한 프로토콜, Privacy를 유지하는 추적 구성, 승인된 소비 차단 중 무엇을 기여로 제시할지 정리합니다.
각 항목이 기존 연구와 어떻게 다른지는 관련연구 비교로 뒷받침해야 합니다.

이어서 보안 분석·모델 검사·구현 실험 중 실제 논문에 포함할 근거를 구분합니다.
\textbf{제안했다는 사실과 안전성·동작·비용을 확인했다는 사실을 별도로 적는 것}이 이 section의 목적입니다.
